\begin{thebibliography}{99}
\small
\setlength{\itemsep}{1pt}
\setlength{\parskip}{0pt}
\bibitem{decreto1469} Gobierno de Colombia (2025). Decreto 1469 de 2025, 29 de diciembre. Fijación del salario mínimo de 2026. \href{https://cancilleria.gov.co/normograma/compilacion/docs/decreto_1469_2025.htm}{Decreto}.
\bibitem{juntaOctubre} Banco de la República (2026). La Junta Directiva decidió mantener inalterada la tasa de interés de política en 12,25\%. Comunicado del 30 de octubre. Expectativas de analistas para diciembre de 2026 de 6,9\%. \href{https://www.banrep.gov.co/es/noticias/junta-directiva/junta-directiva-octubre-2026}{Comunicado}.
\bibitem{geihSalario} DANE (2026). GEIH 2025, diccionario de la variable P6500. Salario mensual monetario del empleo principal antes de descuentos, con propinas y comisiones, sin subsidios ni pagos en especie. \href{https://microdatos.dane.gov.co/index.php/catalog/853/variable/F2/V122?name=P6500}{Diccionario}.
\bibitem{geihTabulados} Centro Javeriano de Competitividad (2026c). Tabulados de salario medio GEIH 2021--2026. Estimaciones de Jorge Mario Orozco Thomas, 2 de noviembre. Enero--agosto de entrevistas en todos los años, con medias, sumas ponderadas, muestras y auditorías. \href{https://github.com/oliverpardo1979/Informe-CJC-SM/tree/main/data/salario_medio_geih}{Datos y metodología}.
\bibitem{ipc} DANE (2026). \textit{IPC. Índices, serie de empalme 2003--2025}. Diciembre de 2025. \href{https://www.dane.gov.co/files/operaciones/IPC/dic2025/anex-IPC-Indices-dic2025.xlsx}{Anexo estadístico}.
\bibitem{productividad} DANE (2026). \textit{Productividad laboral 2005--2025pr}. Cuadros 1 y 2. Publicación del 31 de marzo. \href{https://www.dane.gov.co/files/operaciones/PTF/anex-PTF-ProductividadLaboral-2025.xlsx}{Anexo estadístico}.
\bibitem{salarios} Gobierno de Colombia. Decretos 1785 de 2020, 1724 de 2021, 2613 de 2022, 2292 de 2023 y 1572 de 2024. Montos del salario mínimo de 2021 a 2025. \href{https://cancilleria.gov.co/normograma/compilacion/docs/decreto_2613_2022.htm}{Decreto 2613 de 2022}. \href{https://www.funcionpublica.gov.co/eva/gestornormativo/norma.php?i=257156}{Decreto 1572 de 2024}.
\bibitem{decreto2026} Gobierno de Colombia (2026). \textit{Decreto 159 de 2026}. Artículo 1 y motivación del aumento del salario mínimo. \href{https://cancilleria.gov.co/normograma/compilacion/docs/decreto_0159_2026.htm}{Texto de la norma}.
\bibitem{orozcoGEIH} Centro Javeriano de Competitividad (2026b). \textit{Distribución del salario mensual por edad.} Estimaciones descriptivas con GEIH 2025 y ponderadores de la encuesta.
\bibitem{salarioHistorico} Gobierno de Colombia (2021). \textit{Decreto 1724 de 2021}. Cuadro histórico de salario mínimo mensual e incrementos nominales. \href{https://www.funcionpublica.gov.co/eva/gestornormativo/norma_pdf.php?i=174267}{Texto y cuadro histórico}.
\bibitem{mfmp2026} Ministerio de Hacienda y Crédito Público (2026). \textit{Marco Fiscal de Mediano Plazo 2026}. Junio. Pensiones y cotizaciones, pp. 60--61. Supuesto de deslizamiento, p. 125. Escenarios de riesgo de Colpensiones, pp. 164--166, gráfico 4.14. \href{https://www.minhacienda.gov.co/documents/d/portal/mfmp-2026?download=true}{Documento oficial}.
\bibitem{carf22} Comité Autónomo de la Regla Fiscal (2026). Comunicado 22, 2 de enero. Estimación preliminar y parcial del efecto fiscal del aumento del salario mínimo de 2026, pp. 1--3 y notas 2--4. \href{https://www.carf.gov.co/documents/d/guest/comunicado-22-estimacion-preliminar-efecto-fiscal-de-aumento-de-salario-minimo-2026?download=true}{Comunicado}.
\end{thebibliography}
